\chapter{Thyroid gland}
\index{Thyroid gland}

\medskip
\noindent\textit{Educational reference; always seek professional advice for individual care.}

\section*{Definition and Overview}
The thyroid gland is an endocrine organ that sits low in the front of the neck, wrapped around the upper trachea below the larynx. Shaped like a butterfly, it has two lobes joined by a thin isthmus. Its main products are the iodine-containing hormones thyroxine (T4) and triiodothyronine (T3), which influence metabolic rate, cardiovascular function, thermogenesis, neurological development, and growth. Parafollicular C cells secrete calcitonin, a minor contributor to calcium regulation compared with parathyroid hormone. The gland is controlled by the hypothalamic--pituitary--thyroid axis: hypothalamic TRH stimulates pituitary TSH, which stimulates thyroid hormone synthesis and release; circulating T3 and T4 feedback to keep the system stable. Understanding normal thyroid anatomy and physiology helps patients make sense of blood tests, goitre, nodules, surgery, and radioiodine treatment discussed elsewhere in this book.

\section*{Epidemiology}
Every human depends on continuous thyroid hormone production. Public health relevance appears when iodine intake is inadequate (goitre and cretinism historically), when autoimmune thyroid disease is prevalent, and when neck irradiation or genetic syndromes raise cancer risk. In many countries, mandatory iodine fortification of bread has improved population iodine status, supporting healthy thyroid hormone synthesis, especially in pregnancy. Clinically silent anatomical variants---hemiagenesis, pyramidal lobe, thyroglossal remnants---are not rare on imaging. Ageing changes the distribution of nodules and the interpretation of TSH reference ranges slightly, which is why results are always read in clinical context.

\section*{Aetiology and Pathophysiology}
\begin{itemize}
    \item \textbf{Follicular architecture:} spheres of thyrocytes surround colloid stores of thyroglobulin
    \item \textbf{Iodine trapping:} NIS sodium--iodide symporter concentrates iodide from blood---the basis of radioiodine scans and therapy
    \item \textbf{Hormone synthesis:} iodination of tyrosyl residues on thyroglobulin produces MIT and DIT, coupled to T4 and T3
    \item \textbf{Release and conversion:} mainly T4 is secreted; peripheral deiodinases convert T4 to active T3 in tissues
    \item \textbf{Plasma transport:} TBG and other proteins bind hormones; free fractions are biologically active and usually measured
    \item \textbf{TSH signalling:} drives growth of thyroid tissue as well as hormone output---chronic TSH elevation can enlarge the gland
    \item \textbf{C cells and calcitonin:} embryologically distinct; origin of medullary thyroid cancer when transformed
    \item \textbf{Anatomical relations:} recurrent laryngeal nerves and parathyroid glands neighbour the thyroid---critical in surgery
\end{itemize}

\section*{Clinical Features}
Healthy thyroid function is invisible. When describing the gland clinically, people and clinicians notice:

\begin{itemize}
    \item \textbf{Normal neck:} no visible bulge; gland may be softly palpable in thin people
    \item \textbf{Goitre:} diffuse or nodular enlargement from iodine deficiency, autoimmunity, or multinodular disease
    \item \textbf{Pain or tenderness:} suggests thyroiditis or haemorrhage into a cyst rather than simple Graves disease
    \item \textbf{Compressive symptoms:} dysphagia, orthopnoea, facial plethora on arm elevation (Pemberton sign) with large retrosternal goitres
    \item \textbf{Systemic hormone effects:} hypo- or hyperthyroid symptom clusters when production is abnormal
    \item \textbf{Voice change after surgery:} raises concern for recurrent laryngeal nerve injury
    \item \textbf{Hypocalcaemic symptoms after surgery:} tingling and tetany from parathyroid disturbance
\end{itemize}

\section*{Differential Diagnosis}
Neck lumps are not all thyroid.

\begin{itemize}
    \item \textbf{Thyroid lobe versus lymph node versus salivary pathology:} ultrasound is definitive
    \item \textbf{Thyroglossal duct cyst:} midline moving with tongue protrusion
    \item \textbf{Parathyroid adenoma:} usually impalpable; biochemical hyperparathyroidism
    \item \textbf{Laryngeal and tracheal disease:} primary airway symptoms
    \item \textbf{Diffuse neck swelling from obesity or lipoma:} non-glandular
    \item \textbf{Branchial cleft anomalies:} lateral cystic masses in younger people
\end{itemize}

\section*{When to Seek Medical Care}
See a GP for a new neck lump, progressive swelling, unexplained hormone-type symptoms, or family history of medullary thyroid cancer or MEN syndromes. Mention pregnancy when thyroid tests are planned.

\textbf{Contact your local emergency services for:}
\begin{itemize}
    \item Sudden severe breathing difficulty or stridor with a neck mass
    \item Post-thyroidectomy neck swelling with breathing compromise (possible haematoma)
    \item Severe muscle spasms or seizures after thyroid surgery suggesting critical hypocalcaemia
\end{itemize}

\section*{Diagnosis and Investigations}
\begin{itemize}
    \item \textbf{TSH, free T4, free T3:} functional status of the axis
    \item \textbf{Antibodies:} autoimmune context
    \item \textbf{Ultrasound:} first structural test; no ionising radiation
    \item \textbf{Nuclear medicine thyroid scan and uptake:} functional maps of hot and cold areas
    \item \textbf{CT or MRI:} large retrosternal goitres and complex anatomy---note iodinated CT contrast effects on subsequent radioiodine timing
    \item \textbf{FNA cytology:} nodule diagnosis
    \item \textbf{Laryngoscopy:} voice assessment before and after surgery in selected cases
    \item \textbf{Calcitonin and genetic tests:} medullary cancer pathways
\end{itemize}

\section*{Management}
Management targets disease of the gland rather than ``supporting'' a healthy gland with unproven supplements.

\begin{itemize}
    \item \textbf{Replace missing hormone:} levothyroxine when the gland cannot meet needs
    \item \textbf{Suppress overproduction:} antithyroid drugs, radioiodine ablation, or surgery
    \item \textbf{Observe structural lesions:} many small benign-appearing nodules need only surveillance
    \item \textbf{Operate when indicated:} cancer, suspicious cytology, toxic nodules failing other care, or compression
    \item \textbf{Protect accompanying structures:} parathyroid preservation and nerve monitoring strategies in modern thyroid surgery
    \item \textbf{Pregnancy iodine needs:} routine prenatal care includes adequate iodine via diet and pregnancy multivitamins as advised in many countries---avoid high-dose kelp excess
    \item \textbf{Medicines affecting tests:} biotin can interfere with some assays; list all supplements
\end{itemize}

\section*{Complications}
Disease and treatment of the thyroid can lead to permanent hypothyroidism (often intentional after ablative therapy), hypoparathyroidism, voice changes, neck haematoma, radioiodine side effects (sialadenitis, temporary neck discomfort), and, rarely, advanced compressive or malignant complications. Untreated fetal and neonatal hypothyroidism risks irreversible neurodevelopmental harm---hence newborn screening programmes.

\section*{Prognosis and Outcomes}
A normal thyroid supports lifelong metabolic health without conscious effort. When disease occurs, modern endocrinology usually restores biochemical normality and good quality of life. Surgical and radioiodine pathways are mature and safe in experienced centres. Newborn screening for congenital hypothyroidism is a major public health success, preventing intellectual disability when treatment starts early.

\section*{Prevention and Self-Care}
\begin{itemize}
    \item Use iodised salt and, in pregnancy, follow local advice on prenatal supplements containing iodine
    \item Avoid megadose iodine supplements and unregulated ``thyroid support'' products that may worsen autoimmunity or assays
    \item Do not smoke
    \item Attend newborn screening follow-up if contacted about an abnormal heel-prick result
    \item After thyroidectomy or radioiodine, commit to lifelong hormone replacement and blood-test monitoring
    \item Protect the neck from unnecessary radiation, especially in childhood
    \item Report neck lumps early rather than waiting for compressive symptoms
\end{itemize}

\section*{Sources and Further Reading}
\begin{itemize}
    \item NHS. \textit{Thyroid gland.} https://www.nhs.uk/
    \item Mayo Clinic. \textit{Thyroid gland.} https://www.mayoclinic.org/
    \item MSD Manual. \textit{Thyroid function.} https://www.msdmanuals.com/
    \item MedlinePlus. \textit{Thyroid gland.} https://medlineplus.gov/
    \item Better Health Channel. \textit{Thyroid gland.} https://www.betterhealth.vic.gov.au/
    \item Healthdirect. \textit{Thyroid.} https://www.healthdirect.gov.au/
\end{itemize}
